% --- Archon physics formalization source begin ---
% archon:physics
% archon:covers PhyXMiniProblems/problem_phyx_mini_0686.lean
% archon:source-report reports/phyx_mini/problem_phyx_mini_0686.source.json

\chapter{Physics problem phyx\_mini\_0686}
\label{ch:PhyXMiniProblems_problem_phyx_mini_0686}

\paragraph{Problem source.}
\#\# Physical scenario

A student stands at the edge of a cliff and throws a stone horizontally over the edge with a speed of \textbackslash{}( v\_i = 18.0 \textbackslash{}text\{ m/s\} \textbackslash{}). The cliff is \textbackslash{}( h = 50.0 \textbackslash{}text\{ m\} \textbackslash{}) above a body of water as shown in figure.

\#\# Question

With what speed does the stone land?

\#\# Answer choices

- A: A: \textbackslash{}( 28.4 \textbackslash{}text\{ m/s\} \textbackslash{})

- B: B: \textbackslash{}( 31.6 \textbackslash{}text\{ m/s\} \textbackslash{})

- C: C: \textbackslash{}( 36.1 \textbackslash{}text\{ m/s\} \textbackslash{})

- D: D: \textbackslash{}( 38.3 \textbackslash{}text\{ m/s\} \textbackslash{})

\#\# Figure caption (auxiliary; use the image as primary evidence)

The image depicts a projectile motion diagram involving a person standing at the edge of a cliff. Here's the breakdown:

1. **Person**: A small figure is standing on the left side of the image at the top of a cliff.

2. **Cliff and Ground**: The cliff is on the left, with a flat horizontal surface extending to the right. 

3. **Projectile Path**: A dashed line represents the trajectory of a projectile.

4. **Initial Velocity Vector (\textbackslash{}(\textbackslash{}vec\{v\_i\}\textbackslash{}))**: A red horizontal arrow pointing to the right, indicating the initial velocity of the projectile.

5. **Final Velocity Vector (\textbackslash{}(\textbackslash{}vec\{v\}\textbackslash{}))**: A red arrow pointing diagonally downward to the right, indicating the velocity of the projectile just before impact.

6. **Gravity Vector (\textbackslash{}(\textbackslash{}vec\{g\}\textbackslash{}))**: A purple arrow pointing downward, representing the acceleration due to gravity.

7. **Axes**:
   - \textbackslash{}(x\textbackslash{}) is the horizontal axis.
   - \textbackslash{}(y\textbackslash{}) is the vertical axis.

8. **Height (\textbackslash{}(h\textbackslash{}))**: A vertical double-headed arrow indicates the height of the cliff from which the projectile is launched.

This diagram is typically used to illustrate concepts of kinematics and projectile motion under the influence of gravity.

\#\# Dataset metadata

Kinematics; Physical Model Grounding Reasoning, Multi-Formula Reasoning

\paragraph{Recorded answer/context.}
C: C: \textbackslash{}( 36.1 \textbackslash{}text\{ m/s\} \textbackslash{})

\paragraph{Figure/image path.}
/root/proposal\_for\_physic/science-mango/phyx\_mini\_run/phyx\_data/test\_image/686.png

\paragraph{Formalization target.}
create a compiling Lean file with sorry bodies at `PhyXMiniProblems/problem\_phyx\_mini\_0686.lean`. The Lean declarations must preserve the physical quantities, dimensions or dimensional roles, figure labels, governing-law hypotheses, and final relation expressed by this problem.
Use Mathlib/Physlib names found through LeanExplore where available. If a physics API is missing, introduce faithful local abstractions rather than scalar placeholder aliases.

\begin{theorem}[Physics formalization target]
\label{thm:physics:phyx_mini_0686:target}
\lean{PhyXMiniProblems.ProblemPhyXMini0686.problem_phyx_mini_0686}
\uses{def:physics:phyx-mini-0686:phyxminiproblems-problemphyxmini0686-speedinmeterspersecond, def:physics:phyx-mini-0686:phyxminiproblems-problemphyxmini0686-horizontalcliffprojectilesetup, def:physics:phyx-mini-0686:phyxminiproblems-problemphyxmini0686-matchesprimaryfigure, def:physics:phyx-mini-0686:phyxminiproblems-problemphyxmini0686-matchesproblemdescription, def:physics:phyx-mini-0686:phyxminiproblems-problemphyxmini0686-usesstandardnearearthgravity, def:physics:phyx-mini-0686:phyxminiproblems-problemphyxmini0686-hasphysicalprojectileparameters, def:physics:phyx-mini-0686:phyxminiproblems-problemphyxmini0686-satisfiesuniformgravityprojectilelaws, def:physics:phyx-mini-0686:phyxminiproblems-problemphyxmini0686-exactimpactspeedinmeterspersecond, def:physics:phyx-mini-0686:phyxminiproblems-problemphyxmini0686-matchesdisplayedimpactspeed, def:physics:phyx-mini-0686:phyxminiproblems-problemphyxmini0686-isuniqueclosestimpactspeedchoice}
The assigned autoformalize agent should translate this physics problem into Lean declarations in the covered file, with theorem and lemma proof bodies written as `by sorry`.
\end{theorem}
\begin{proof}
This is an autoformalization task, not a proof task. Produce statements that can later be proved by the physics prover without weakening the physical model.
\end{proof}

% --- Archon declaration topology begin ---
\section{Declaration topology}
\begin{definition}[velocity Dimension]
  \label{def:physics:phyx-mini-0686:phyxminiproblems-problemphyxmini0686-velocitydimension}
  \lean{PhyXMiniProblems.ProblemPhyXMini0686.velocityDimension}
  The physical dimension of velocity, L T⁻¹
\end{definition}
\begin{proof}
  This is the declared model definition of velocity Dimension.
\end{proof}
\begin{definition}[acceleration Dimension]
  \label{def:physics:phyx-mini-0686:phyxminiproblems-problemphyxmini0686-accelerationdimension}
  \lean{PhyXMiniProblems.ProblemPhyXMini0686.accelerationDimension}
  The physical dimension of acceleration, L T⁻²
\end{definition}
\begin{proof}
  This is the declared model definition of acceleration Dimension.
\end{proof}
\begin{definition}[Length Quantity]
  \label{def:physics:phyx-mini-0686:phyxminiproblems-problemphyxmini0686-lengthquantity}
  \lean{PhyXMiniProblems.ProblemPhyXMini0686.LengthQuantity}
  A nonnegative physical length
\end{definition}
\begin{proof}
  This is the declared model definition of Length Quantity.
\end{proof}
\begin{definition}[Time Quantity]
  \label{def:physics:phyx-mini-0686:phyxminiproblems-problemphyxmini0686-timequantity}
  \lean{PhyXMiniProblems.ProblemPhyXMini0686.TimeQuantity}
  A nonnegative physical duration
\end{definition}
\begin{proof}
  This is the declared model definition of Time Quantity.
\end{proof}
\begin{definition}[Planar Position Quantity]
  \label{def:physics:phyx-mini-0686:phyxminiproblems-problemphyxmini0686-planarpositionquantity}
  \lean{PhyXMiniProblems.ProblemPhyXMini0686.PlanarPositionQuantity}
  A signed position vector in the plane
\end{definition}
\begin{proof}
  This is the declared model definition of Planar Position Quantity.
\end{proof}
\begin{definition}[Planar Velocity Quantity]
  \label{def:physics:phyx-mini-0686:phyxminiproblems-problemphyxmini0686-planarvelocityquantity}
  \lean{PhyXMiniProblems.ProblemPhyXMini0686.PlanarVelocityQuantity}
  \uses{def:physics:phyx-mini-0686:phyxminiproblems-problemphyxmini0686-velocitydimension}
  A signed velocity vector in the plane
\end{definition}
\begin{proof}
  This is the declared model definition of Planar Velocity Quantity.
\end{proof}
\begin{definition}[Planar Acceleration Quantity]
  \label{def:physics:phyx-mini-0686:phyxminiproblems-problemphyxmini0686-planaraccelerationquantity}
  \lean{PhyXMiniProblems.ProblemPhyXMini0686.PlanarAccelerationQuantity}
  \uses{def:physics:phyx-mini-0686:phyxminiproblems-problemphyxmini0686-accelerationdimension}
  A signed acceleration vector in the plane
\end{definition}
\begin{proof}
  This is the declared model definition of Planar Acceleration Quantity.
\end{proof}
\begin{definition}[x Axis]
  \label{def:physics:phyx-mini-0686:phyxminiproblems-problemphyxmini0686-xaxis}
  \lean{PhyXMiniProblems.ProblemPhyXMini0686.xAxis}
  The coordinate index carrying the figure's x label
\end{definition}
\begin{proof}
  This is the declared model definition of x Axis.
\end{proof}
\begin{definition}[y Axis]
  \label{def:physics:phyx-mini-0686:phyxminiproblems-problemphyxmini0686-yaxis}
  \lean{PhyXMiniProblems.ProblemPhyXMini0686.yAxis}
  The coordinate index carrying the figure's y label
\end{definition}
\begin{proof}
  This is the declared model definition of y Axis.
\end{proof}
\begin{definition}[length In Meters]
  \label{def:physics:phyx-mini-0686:phyxminiproblems-problemphyxmini0686-lengthinmeters}
  \lean{PhyXMiniProblems.ProblemPhyXMini0686.lengthInMeters}
  \uses{def:physics:phyx-mini-0686:phyxminiproblems-problemphyxmini0686-lengthquantity}
  Metre readout of a nonnegative physical length
\end{definition}
\begin{proof}
  This is the declared model definition of length In Meters.
\end{proof}
\begin{definition}[time In Seconds]
  \label{def:physics:phyx-mini-0686:phyxminiproblems-problemphyxmini0686-timeinseconds}
  \lean{PhyXMiniProblems.ProblemPhyXMini0686.timeInSeconds}
  \uses{def:physics:phyx-mini-0686:phyxminiproblems-problemphyxmini0686-timequantity}
  Second readout of a nonnegative physical duration
\end{definition}
\begin{proof}
  This is the declared model definition of time In Seconds.
\end{proof}
\begin{definition}[position In Meters]
  \label{def:physics:phyx-mini-0686:phyxminiproblems-problemphyxmini0686-positioninmeters}
  \lean{PhyXMiniProblems.ProblemPhyXMini0686.positionInMeters}
  \uses{def:physics:phyx-mini-0686:phyxminiproblems-problemphyxmini0686-planarpositionquantity}
  Coherent-SI position vector, whose coordinates are measured in metres
\end{definition}
\begin{proof}
  This is the declared model definition of position In Meters.
\end{proof}
\begin{definition}[velocity In Meters Per Second]
  \label{def:physics:phyx-mini-0686:phyxminiproblems-problemphyxmini0686-velocityinmeterspersecond}
  \lean{PhyXMiniProblems.ProblemPhyXMini0686.velocityInMetersPerSecond}
  \uses{def:physics:phyx-mini-0686:phyxminiproblems-problemphyxmini0686-planarvelocityquantity}
  Coherent-SI velocity vector, in metres per second
\end{definition}
\begin{proof}
  This is the declared model definition of velocity In Meters Per Second.
\end{proof}
\begin{definition}[acceleration In Meters Per Second Squared]
  \label{def:physics:phyx-mini-0686:phyxminiproblems-problemphyxmini0686-accelerationinmeterspersecondsquared}
  \lean{PhyXMiniProblems.ProblemPhyXMini0686.accelerationInMetersPerSecondSquared}
  \uses{def:physics:phyx-mini-0686:phyxminiproblems-problemphyxmini0686-planaraccelerationquantity}
  Coherent-SI acceleration vector, in metres per second squared
\end{definition}
\begin{proof}
  This is the declared model definition of acceleration In Meters Per Second Squared.
\end{proof}
\begin{definition}[speed In Meters Per Second]
  \label{def:physics:phyx-mini-0686:phyxminiproblems-problemphyxmini0686-speedinmeterspersecond}
  \lean{PhyXMiniProblems.ProblemPhyXMini0686.speedInMetersPerSecond}
  Metres-per-second readout of a nonnegative physical speed
\end{definition}
\begin{proof}
  This is the declared model definition of speed In Meters Per Second.
\end{proof}
\begin{definition}[Projectile Kind]
  \label{def:physics:phyx-mini-0686:phyxminiproblems-problemphyxmini0686-projectilekind}
  \lean{PhyXMiniProblems.ProblemPhyXMini0686.ProjectileKind}
  The kind of projectile named in the prose
\end{definition}
\begin{proof}
  This is the declared model definition of Projectile Kind.
\end{proof}
\begin{definition}[Arrow Direction]
  \label{def:physics:phyx-mini-0686:phyxminiproblems-problemphyxmini0686-arrowdirection}
  \lean{PhyXMiniProblems.ProblemPhyXMini0686.ArrowDirection}
  Qualitative arrow directions visible in the primary image
\end{definition}
\begin{proof}
  This is the declared model definition of Arrow Direction.
\end{proof}
\begin{definition}[Figure Element]
  \label{def:physics:phyx-mini-0686:phyxminiproblems-problemphyxmini0686-figureelement}
  \lean{PhyXMiniProblems.ProblemPhyXMini0686.FigureElement}
  Objects and distinguished points visible in the cliff diagram
\end{definition}
\begin{proof}
  This is the declared model definition of Figure Element.
\end{proof}
\begin{definition}[Figure Quantity Label]
  \label{def:physics:phyx-mini-0686:phyxminiproblems-problemphyxmini0686-figurequantitylabel}
  \lean{PhyXMiniProblems.ProblemPhyXMini0686.FigureQuantityLabel}
  Vector and distance labels printed in the primary image
\end{definition}
\begin{proof}
  This is the declared model definition of Figure Quantity Label.
\end{proof}
\begin{definition}[Cliff Projectile Figure]
  \label{def:physics:phyx-mini-0686:phyxminiproblems-problemphyxmini0686-cliffprojectilefigure}
  \lean{PhyXMiniProblems.ProblemPhyXMini0686.CliffProjectileFigure}
  \uses{def:physics:phyx-mini-0686:phyxminiproblems-problemphyxmini0686-arrowdirection, def:physics:phyx-mini-0686:phyxminiproblems-problemphyxmini0686-figureelement, def:physics:phyx-mini-0686:phyxminiproblems-problemphyxmini0686-figurequantitylabel}
  Literal qualitative and schematic-coordinate content of image 686.png
\end{definition}
\begin{proof}
  This is the declared model definition of Cliff Projectile Figure.
\end{proof}
\begin{definition}[Horizontal Cliff Projectile Setup]
  \label{def:physics:phyx-mini-0686:phyxminiproblems-problemphyxmini0686-horizontalcliffprojectilesetup}
  \lean{PhyXMiniProblems.ProblemPhyXMini0686.HorizontalCliffProjectileSetup}
  \uses{def:physics:phyx-mini-0686:phyxminiproblems-problemphyxmini0686-lengthquantity, def:physics:phyx-mini-0686:phyxminiproblems-problemphyxmini0686-timequantity, def:physics:phyx-mini-0686:phyxminiproblems-problemphyxmini0686-planarpositionquantity, def:physics:phyx-mini-0686:phyxminiproblems-problemphyxmini0686-planarvelocityquantity, def:physics:phyx-mini-0686:phyxminiproblems-problemphyxmini0686-planaraccelerationquantity, def:physics:phyx-mini-0686:phyxminiproblems-problemphyxmini0686-projectilekind, def:physics:phyx-mini-0686:phyxminiproblems-problemphyxmini0686-cliffprojectilefigure}
  Independent physical quantities in the launch-and-impact experiment. In particular, impactSpeed is an independent field. It is related to the impact velocity only by the governing-law premise below and is not defined from the recorded numerical answer
\end{definition}
\begin{proof}
  This is the declared model definition of Horizontal Cliff Projectile Setup.
\end{proof}
\begin{definition}[Matches Primary Figure]
  \label{def:physics:phyx-mini-0686:phyxminiproblems-problemphyxmini0686-matchesprimaryfigure}
  \lean{PhyXMiniProblems.ProblemPhyXMini0686.MatchesPrimaryFigure}
  \uses{def:physics:phyx-mini-0686:phyxminiproblems-problemphyxmini0686-horizontalcliffprojectilesetup}
  Primary-image evidence: the student is at the cliff-top launch point, the water and impact point are below it, the impact is to the right, and the three vector arrows have the directions drawn in the raster image
\end{definition}
\begin{proof}
  This is the declared model definition of Matches Primary Figure.
\end{proof}
\begin{definition}[Matches Problem Description]
  \label{def:physics:phyx-mini-0686:phyxminiproblems-problemphyxmini0686-matchesproblemdescription}
  \lean{PhyXMiniProblems.ProblemPhyXMini0686.MatchesProblemDescription}
  \uses{def:physics:phyx-mini-0686:phyxminiproblems-problemphyxmini0686-xaxis, def:physics:phyx-mini-0686:phyxminiproblems-problemphyxmini0686-yaxis, def:physics:phyx-mini-0686:phyxminiproblems-problemphyxmini0686-lengthinmeters, def:physics:phyx-mini-0686:phyxminiproblems-problemphyxmini0686-positioninmeters, def:physics:phyx-mini-0686:phyxminiproblems-problemphyxmini0686-velocityinmeterspersecond, def:physics:phyx-mini-0686:phyxminiproblems-problemphyxmini0686-horizontalcliffprojectilesetup}
  Numerical and geometric data stated in the prose. The two component equations express that the 18.0 m/s launch is horizontal in the figure's coordinates. The final speed and every answer choice remain absent from these premises
\end{definition}
\begin{proof}
  This is the declared model definition of Matches Problem Description.
\end{proof}
\begin{definition}[Uses Standard Near Earth Gravity]
  \label{def:physics:phyx-mini-0686:phyxminiproblems-problemphyxmini0686-usesstandardnearearthgravity}
  \lean{PhyXMiniProblems.ProblemPhyXMini0686.UsesStandardNearEarthGravity}
  \uses{def:physics:phyx-mini-0686:phyxminiproblems-problemphyxmini0686-xaxis, def:physics:phyx-mini-0686:phyxminiproblems-problemphyxmini0686-yaxis, def:physics:phyx-mini-0686:phyxminiproblems-problemphyxmini0686-accelerationinmeterspersecondsquared, def:physics:phyx-mini-0686:phyxminiproblems-problemphyxmini0686-horizontalcliffprojectilesetup}
  Standard near-Earth gravitational acceleration in the drawn coordinates
\end{definition}
\begin{proof}
  This is the declared model definition of Uses Standard Near Earth Gravity.
\end{proof}
\begin{definition}[Has Physical Projectile Parameters]
  \label{def:physics:phyx-mini-0686:phyxminiproblems-problemphyxmini0686-hasphysicalprojectileparameters}
  \lean{PhyXMiniProblems.ProblemPhyXMini0686.HasPhysicalProjectileParameters}
  \uses{def:physics:phyx-mini-0686:phyxminiproblems-problemphyxmini0686-lengthinmeters, def:physics:phyx-mini-0686:phyxminiproblems-problemphyxmini0686-timeinseconds, def:physics:phyx-mini-0686:phyxminiproblems-problemphyxmini0686-speedinmeterspersecond, def:physics:phyx-mini-0686:phyxminiproblems-problemphyxmini0686-horizontalcliffprojectilesetup}
  Positivity conditions selecting a physically meaningful flight
\end{definition}
\begin{proof}
  This is the declared model definition of Has Physical Projectile Parameters.
\end{proof}
\begin{definition}[Satisfies Uniform Gravity Projectile Laws]
  \label{def:physics:phyx-mini-0686:phyxminiproblems-problemphyxmini0686-satisfiesuniformgravityprojectilelaws}
  \lean{PhyXMiniProblems.ProblemPhyXMini0686.SatisfiesUniformGravityProjectileLaws}
  \uses{def:physics:phyx-mini-0686:phyxminiproblems-problemphyxmini0686-timeinseconds, def:physics:phyx-mini-0686:phyxminiproblems-problemphyxmini0686-positioninmeters, def:physics:phyx-mini-0686:phyxminiproblems-problemphyxmini0686-velocityinmeterspersecond, def:physics:phyx-mini-0686:phyxminiproblems-problemphyxmini0686-accelerationinmeterspersecondsquared, def:physics:phyx-mini-0686:phyxminiproblems-problemphyxmini0686-speedinmeterspersecond, def:physics:phyx-mini-0686:phyxminiproblems-problemphyxmini0686-horizontalcliffprojectilesetup}
  This structure records the Satisfies Uniform Gravity Projectile Laws component of the problem-specific physical model
\end{definition}
\begin{proof}
  This is the declared model definition of Satisfies Uniform Gravity Projectile Laws.
\end{proof}
\begin{definition}[exact Impact Speed In Meters Per Second]
  \label{def:physics:phyx-mini-0686:phyxminiproblems-problemphyxmini0686-exactimpactspeedinmeterspersecond}
  \lean{PhyXMiniProblems.ProblemPhyXMini0686.exactImpactSpeedInMetersPerSecond}
  This def records the exact Impact Speed In Meters Per Second component of the problem-specific physical model
\end{definition}
\begin{proof}
  This is the declared model definition of exact Impact Speed In Meters Per Second.
\end{proof}
\begin{definition}[Answer Choice]
  \label{def:physics:phyx-mini-0686:phyxminiproblems-problemphyxmini0686-answerchoice}
  \lean{PhyXMiniProblems.ProblemPhyXMini0686.AnswerChoice}
  Labels of the four speed choices printed with the problem
\end{definition}
\begin{proof}
  This is the declared model definition of Answer Choice.
\end{proof}
\begin{definition}[speed Meters Per Second]
  \label{def:physics:phyx-mini-0686:phyxminiproblems-problemphyxmini0686-answerchoice-speedmeterspersecond}
  \lean{PhyXMiniProblems.ProblemPhyXMini0686.AnswerChoice.speedMetersPerSecond}
  \uses{def:physics:phyx-mini-0686:phyxminiproblems-problemphyxmini0686-answerchoice}
  Speed in metres per second printed beside each displayed choice
\end{definition}
\begin{proof}
  This is the declared model definition of speed Meters Per Second.
\end{proof}
\begin{definition}[recorded Answer Choice]
  \label{def:physics:phyx-mini-0686:phyxminiproblems-problemphyxmini0686-recordedanswerchoice}
  \lean{PhyXMiniProblems.ProblemPhyXMini0686.recordedAnswerChoice}
  \uses{def:physics:phyx-mini-0686:phyxminiproblems-problemphyxmini0686-answerchoice}
  The answer label recorded by the supplied dataset metadata
\end{definition}
\begin{proof}
  This is the declared model definition of recorded Answer Choice.
\end{proof}
\begin{definition}[answer Choice Error Meters Per Second]
  \label{def:physics:phyx-mini-0686:phyxminiproblems-problemphyxmini0686-answerchoiceerrormeterspersecond}
  \lean{PhyXMiniProblems.ProblemPhyXMini0686.answerChoiceErrorMetersPerSecond}
  \uses{def:physics:phyx-mini-0686:phyxminiproblems-problemphyxmini0686-speedinmeterspersecond, def:physics:phyx-mini-0686:phyxminiproblems-problemphyxmini0686-horizontalcliffprojectilesetup, def:physics:phyx-mini-0686:phyxminiproblems-problemphyxmini0686-answerchoice}
  Absolute error between the modeled impact speed and a displayed choice
\end{definition}
\begin{proof}
  This is the declared model definition of answer Choice Error Meters Per Second.
\end{proof}
\begin{definition}[Matches Displayed Impact Speed]
  \label{def:physics:phyx-mini-0686:phyxminiproblems-problemphyxmini0686-matchesdisplayedimpactspeed}
  \lean{PhyXMiniProblems.ProblemPhyXMini0686.MatchesDisplayedImpactSpeed}
  \uses{def:physics:phyx-mini-0686:phyxminiproblems-problemphyxmini0686-horizontalcliffprojectilesetup, def:physics:phyx-mini-0686:phyxminiproblems-problemphyxmini0686-answerchoice, def:physics:phyx-mini-0686:phyxminiproblems-problemphyxmini0686-answerchoiceerrormeterspersecond}
  Agreement with a speed displayed to the nearest tenth of a metre per second
\end{definition}
\begin{proof}
  This is the declared model definition of Matches Displayed Impact Speed.
\end{proof}
\begin{definition}[Is Unique Closest Impact Speed Choice]
  \label{def:physics:phyx-mini-0686:phyxminiproblems-problemphyxmini0686-isuniqueclosestimpactspeedchoice}
  \lean{PhyXMiniProblems.ProblemPhyXMini0686.IsUniqueClosestImpactSpeedChoice}
  \uses{def:physics:phyx-mini-0686:phyxminiproblems-problemphyxmini0686-horizontalcliffprojectilesetup, def:physics:phyx-mini-0686:phyxminiproblems-problemphyxmini0686-answerchoice, def:physics:phyx-mini-0686:phyxminiproblems-problemphyxmini0686-answerchoiceerrormeterspersecond}
  A displayed choice is strictly closer to the modeled speed than all others
\end{definition}
\begin{proof}
  This is the declared model definition of Is Unique Closest Impact Speed Choice.
\end{proof}
\begin{lemma}[impact Speed eq exact Expression]
  \label{lem:physics:phyx-mini-0686:phyxminiproblems-problemphyxmini0686-impactspeed-eq-exactexpression}
  \lean{PhyXMiniProblems.ProblemPhyXMini0686.impactSpeed_eq_exactExpression}
  \uses{def:physics:phyx-mini-0686:phyxminiproblems-problemphyxmini0686-speedinmeterspersecond, def:physics:phyx-mini-0686:phyxminiproblems-problemphyxmini0686-horizontalcliffprojectilesetup, def:physics:phyx-mini-0686:phyxminiproblems-problemphyxmini0686-matchesprimaryfigure, def:physics:phyx-mini-0686:phyxminiproblems-problemphyxmini0686-matchesproblemdescription, def:physics:phyx-mini-0686:phyxminiproblems-problemphyxmini0686-usesstandardnearearthgravity, def:physics:phyx-mini-0686:phyxminiproblems-problemphyxmini0686-hasphysicalprojectileparameters, def:physics:phyx-mini-0686:phyxminiproblems-problemphyxmini0686-satisfiesuniformgravityprojectilelaws, def:physics:phyx-mini-0686:phyxminiproblems-problemphyxmini0686-exactimpactspeedinmeterspersecond}
  The governing laws determine the exact, unrounded impact-speed readout
\end{lemma}
\begin{proof}
  The conclusion follows from the governing relations and figure readouts named in the statement of impact Speed eq exact Expression.
\end{proof}
\begin{lemma}[exact Impact Speed selects choice C]
  \label{lem:physics:phyx-mini-0686:phyxminiproblems-problemphyxmini0686-exactimpactspeed-selects-choice-c}
  \lean{PhyXMiniProblems.ProblemPhyXMini0686.exactImpactSpeed_selects_choice_C}
  \uses{def:physics:phyx-mini-0686:phyxminiproblems-problemphyxmini0686-speedinmeterspersecond, def:physics:phyx-mini-0686:phyxminiproblems-problemphyxmini0686-horizontalcliffprojectilesetup, def:physics:phyx-mini-0686:phyxminiproblems-problemphyxmini0686-exactimpactspeedinmeterspersecond, def:physics:phyx-mini-0686:phyxminiproblems-problemphyxmini0686-matchesdisplayedimpactspeed, def:physics:phyx-mini-0686:phyxminiproblems-problemphyxmini0686-isuniqueclosestimpactspeedchoice}
  The exact expression is approximately 36.11 m/s; hence it rounds to the displayed value 36.1 m/s and is closer to choice C than to the other choices
\end{lemma}
\begin{proof}
  The conclusion follows from the governing relations and figure readouts named in the statement of exact Impact Speed selects choice C.
\end{proof}
% --- Archon declaration topology end ---
% --- Archon physics formalization source end ---
